% =====================================================================
%  Appendix A -- the source document's own section map.
%  Source: tier1.md:60-136 (the Contents table and the <details> "Full
%  section list"), reproduced so that the source's own Part I--VIII
%  architecture survives the re-architecture into layers.
% =====================================================================

\section{The source document's section map}
\label{app:sectionmap}

This report re-architects \code{study/tier1.md} into epistemic layers
(\S\ref{sec:status}). The source's own organisation --- node by node, Part~I to
Part~VIII --- is reproduced here so that nothing about it is lost, and so that a
reader who knows the source document can navigate this one.
Appendix~\ref{app:concordance} is the heading-by-heading map between them.

\subsection{The source's Contents table}

\begin{center}
\begin{tabularx}{\textwidth}{@{}L{3.0cm} Y@{}}
\toprule
\textbf{Part} & \textbf{Covers} \\
\midrule
\textbf{I} --- Rate theory &
  \sigv{} from cross-sections, Coulomb barrier \& Gamow peak, resonances \&
  Hauser--Feshbach, photodisintegration, what the rates \emph{do} in Si
  burning \\
\textbf{II} (S2) --- REACLIB &
  The seven-coefficient form derived, sets\,/\,chapters\,/\,labels,
  $\lambda \to$ gross flux $R_j$, the compiled evaluator \\
\textbf{III} (S3) --- Detailed balance &
  The reverse\,/\,forward ratio derived in full, the phase-space factor per
  $|\Delta N|$, partition functions, the v-flag dissection, $\kappa$ as the DB
  test, gh-575 \\
\textbf{IV} (S4) --- Screening &
  Debye--H\"uckel derived, $\Gamma$ regimes, \code{chugunov\_2007}
  term-by-term, the per-reaction pairing and its $\kappa$ consequence \\
\textbf{V} (S5) --- The weak sector &
  Fermi golden rule $\to$ $ft$ $\to$ allowed rates, Gamow--Teller, EC in a
  degenerate plasma, the $(T, \rhoye)$ tables, table families, the network
  asymmetry \\
\textbf{VI} (S6) --- Reconciliation &
  Canonical keys, membership vs values, dispositions, the measured comparison,
  the Fortran probe, what it licenses \\
\textbf{VII} --- Executable policy &
  \code{fluxes/guards.py}: the four Tier-1 gates as code that refuses to run \\
\textbf{VIII} --- Open prerequisites &
  Standing cross-tier register. Two audit passes: the physics the code
  implements (\S VIII.C, ten gaps) and the project's own spec documents vs the
  study map (\S VIII.E --- an entire missing \textbf{Tier 4}: the model, its
  conditioning, its training) \\
\bottomrule
\end{tabularx}
\end{center}

\subsection{The source's full section list}

Reproduced verbatim from the source's collapsed ``Full section list'' block.

\begin{srclisting}
1.1  From cross-section to <sigma v>  1.1.1 rate per volume . 1.1.2 double counting
                                      1.1.3 the molar convention . 1.1.4 code contract
1.2  Coulomb barrier & Gamow peak     1.2.1 tunnelling . 1.2.2 the S-factor
                                      1.2.3 saddle point -> E0, Delta
                                      1.2.4 numbers in the box
                                      1.2.5 neutrons have no Gamow peak
                                      1.2.6 the 1/v law . the reciprocity theorem
1.3  Resonances & Hauser-Feshbach     narrow-resonance rate . the HF regime
                                      1.3.3 where the uncertainty actually is
                                      1.3.4 Gamma/D, level densities, Wigner limit
                                      1.3.5 (!) ground-state vs stellar rates (SEF)
1.4  Photodisintegration              the Planck tail . the Q/kT exponential
                                      worked crossover Si28 <-> S32
                                      why QSE turns on at ~3 GK
                                      1.4.4 LTE for photons . 3-alpha & the Hoyle state
1.5  What the rates do                mapping Part I onto the Tier-0 topology
1.6  self-check

II.1 The seven-coefficient form    II.5  Code: the coefficient tensor
II.2 Sets, chapters, labels        II.6  Oracle: the 5e-12 gate
II.3 One rate dissected            II.6b lambda's 75-decade range; over/underflow
II.4 lambda -> R_j                 II.6c multiple sets; fit validity range
                                   II.7  A fit library, not a theory
                                   II.8  self-check

III.1 DB from chemical equilibrium III.7  gh-575 derived from the same algebra
III.2 The phase-space factor       III.8  Code: db_reverses + the pf gate
III.3 Partition functions          III.9  Oracles
      (!) the G vs (2J+1)G trap    III.10 The two kappa conventions
      spline extrapolation         III.11 Coulomb corrections to NSE
III.4 Three ways to build a reverse III.12 Q(T) derived + measured
III.5 The v-flag dissected         III.13 self-check
III.6 kappa as the numerical DB test

IV.1 Why screening exists          IV.6  Per-reaction pairs
IV.2 Debye-Huckel derived          IV.7  The screened-kappa offset derived
IV.3 Strong screening & Gamma      IV.8  Code + oracle
      the free-energy derivation   IV.9  Screening of the WEAK sector (MU/DQ/VS)
IV.4 Where the box sits            IV.10 self-check
IV.5 chugunov_2007 term by term

V.1 Why weak is categorical        V.6  The three channels & the ledger
V.2 Golden rule -> ft              V.7  The table families
V.3 Gamow-Teller strength          V.8  Neutrino losses
V.4 EC in a degenerate plasma      V.9  Code
      V.4b Pauli blocking = the ratchet   V.10 The mesa_80/151 asymmetry
V.5 The (T, rho*Ye) table          V.11 self-check
      V.5b interpolation error derived

VI.1 Why this node exists          VI.6  Appendix-B as a measurement
VI.2 The canonical key             VI.7  The screening determination
VI.3 Membership vs values          VI.8  The Fortran probe
VI.4 The disposition algebra       VI.8b probe modes & the grid
VI.5 The measured comparison       VI.8c reading the statistics honestly
                                   VI.9  What it licenses . VI.10 self-check

VIII.A deferred by design          VIII.C (!) ten genuine gaps (physics pass)
VIII.B closed in this revision     VIII.D the common-mode habit
                                   VIII.E (!!) the missing TIER 4 (spec pass)
                                   VIII.F priorities across both passes
\end{srclisting}

\begin{aside}
\textbf{Glyph substitutions in the listing above.} The block is set in a
monospaced face, so a handful of source glyphs were transliterated rather than
risk a missing glyph: \sigv{} $\to$ \code{<sigma v>}, $\Tn \to$ \code{T9},
$E_0 \to$ \code{E0}, \rhoye{} $\to$ \code{rho*Ye}, $\rightleftharpoons\ \to$
\code{<->}, $\to\ \to$ \code{->}, the interpunct $\to$ \code{.}, Greek letters
to their names, isotope superscripts to \code{Si28}\,/\,\code{S32}, and the
warning glyph to \code{(!)} (doubled as \code{(!!)} where the source doubles
it). Nothing else was changed: the line breaks, the two-column layout, and the
ordering are the source's.
\end{aside}
